\bookchapter{Capítulo 56: Pulir y publicar: pruebas, depuración y salida a Google Play}

\chapterintro{Ya tienes la App, ¿y ahora qué? Este capítulo responde a las tres tareas que separan “que funcione” de “que esté lista para entregar”: cómo depurar con eficacia para encontrar bugs, cómo escribir pruebas que eviten regresiones y cómo subir la App a Google Play para que todo el mundo la descargue. Publicarla no es solo un acto ceremonial: es la línea que separa una obra de un producto.}

\booksection{Las tres herramientas de depuración: Logcat, puntos de interrupción e inspector de UI}

\begin{center}
\begin{tabularx}{\linewidth}{|X|X|X|}
\hline
\textbf{Herramienta} & \textbf{Para qué sirve} & \textbf{Cómo se usa} \tabularnewline
\hline
Logcat & Ver los registros del sistema o los tuyos propios & En el código escribe Log.d("TAG", "valor=\$x") y filtra buscando el TAG en el panel \tabularnewline
Depuración con puntos de interrupción & Detener el programa en una línea y ver los valores de las variables & Haz clic junto al número de línea para que salga un punto rojo, ejecuta en modo Debug y avanza paso a paso \tabularnewline
Layout Inspector & Ver la estructura de la interfaz en ejecución & Menú Tools Layout Inspector, para ver el árbol de componentes en tiempo real \tabularnewline
\hline
\end{tabularx}
\end{center}

\begin{center}\small\color{gray}Las tres herramientas de depuración\end{center}

\begin{tcbinfo}{Consejo: ¿cómo leer un error?}
Tras un cierre inesperado, en Logcat aparece un stacktrace en rojo. Busca de abajo hacia arriba la primera línea que empiece por el nombre de tu paquete (por ejemplo, com.example.yourapp); esa línea es casi siempre la escena del crimen. Mira primero el “tipo de Exception” (¿puntero nulo?, ¿conversión de tipos?) y luego el “at ... tuarchivo: línea”. Este truco resuelve el 80\% de los cierres inesperados.

\end{tcbinfo}

\booksection{Las pruebas: un seguro para tu código}

Las pruebas manuales no son de fiar: arreglas una función y otra se rompe en silencio. Las pruebas son a la vez el remedio arrepentido y el chaleco antibalas. Las pruebas en Android se dividen en tres niveles; si empiezas, domina primero el primero:

\begin{center}
\begin{tabularx}{\linewidth}{|X|X|X|X|}
\hline
\textbf{Nivel} & \textbf{Qué prueba} & \textbf{Dónde va} & \textbf{Velocidad} \tabularnewline
\hline
Pruebas unitarias & Lógica pura (cálculos, validaciones) & app/src/test/ & Segundos \tabularnewline
Pruebas instrumentadas & Comportamiento de la interfaz en un dispositivo real & app/src/androidTest/ & Minutos \tabularnewline
Pruebas manuales & Percepción y experiencia de una persona real & —— & Siempre necesarias \tabularnewline
\hline
\end{tabularx}
\end{center}

\begin{center}\small\color{gray}Los tres niveles de pruebas\end{center}

\begin{lstlisting}[language=Kotlin]
// extrae a una función pura la lógica de "convertir el texto del importe en número"
fun parseAmount(input: String): Double? {
    val v = input.trim().toDoubleOrNull()
    return if (v != null && v > 0) v else null
}

// pruébala (JUnit4)
@Test
fun emptyInputReturnsNull() {
    assertEquals(null, parseAmount(""))
}
@Test
fun negativeReturnsNull() {
    assertEquals(null, parseAmount("-5"))
}
@Test
fun normalValueParses() {
    assertEquals(12.5, parseAmount("12.5"))
}
\end{lstlisting}

La validación de importes y los cálculos estadísticos de tu app de gastos deberían estar cubiertos por pruebas unitarias. Para ejecutarlas: clic derecho en la carpeta test Run. Si todo sale en verde, puedes cambiar el código con tranquilidad.

\booksection{Firma y empaquetado: ponerle el “sello” a tu App}

Android exige que todos los paquetes de instalación estén firmados: es como un sello oficial que certifica “este paquete sale de mí y no ha sido alterado”. Al depurar, Android Studio usa automáticamente la firma debug; para publicar necesitas tu propia firma release. El proceso:

\begin{enumerate}
  \item Menú Build Generate Signed Bundle / APK;
  \item Crea el archivo keystore (almacén de claves): define una contraseña y guárdala bien;
  \item Elige release y el formato AAB, y la compilación queda lista.
\end{enumerate}

\textbf{Advertencia: perder el keystore = perder la App}\par

Para actualizar la App después de publicarla debes usar la misma firma; si no, la tienda rechazará la actualización. Si pierdes el archivo keystore o la contraseña, solo te queda publicar una app nueva desde cero, y los usuarios antiguos se quedan sin poder actualizar. Haz siempre una copia de seguridad del keystore y de la contraseña en un lugar seguro (nube cifrada + un bloc de notas): es una de las lecciones más caras para un desarrollador Android.

\booksection{Todo el proceso para publicar en Google Play}

Registrar cuenta de desarrollador

Preparar material de la tienda

Subir el AAB

Rellenar el formulario de seguridad de datos

Revisión (de horas a días)

Publicar

\begin{center}\small\color{gray}Figura 6: los seis pasos para publicar en Google Play\end{center}

\begin{itemize}
  \item Registro: una cuota única de 25 USD y una cuenta de Google que funcione;
  \item Material: icono de la app (512×512), capturas de pantalla (al menos 2), descripción de funciones y enlace a la política de privacidad (obligatorio si recopilas datos);
  \item Formulario de seguridad de datos: declara con honestidad qué datos recopilas; una app de gastos sin conexión puede declarar "cero recopilación";
  \item Revisión: la primera publicación de una cuenta nueva suele pasar por revisión humana, que tarda de horas a días; si la rechazan, corrige según los comentarios.
\end{itemize}

Nota para usuarios de China: Google Play requiere un entorno de red especial y la distribución en China continental suele hacerse a través de tiendas nacionales (Huawei, Xiaomi, MyApp de Tencent, etc.). El proceso para publicar en esas tiendas es parecido, pero cada una tiene sus propios criterios de revisión; el mismo AAB/APK puede distribuirse en varias tiendas. Para practicar, con lograr publicar en una ya ganas.

\booksection{Cómo crear el material de la tienda: la primera impresión son las descargas}

La página de la tienda es el escaparate que ven los usuarios al pasear por ella. Tres elementos deciden la primera impresión:

\begin{itemize}
  \item Icono: 512×512; no uses una captura de pantalla como parche. Si no tienes ojo de diseñador, genera con una herramienta de IA una figura geométrica sencilla; solo vale si se reconoce de un vistazo incluso en tamaño reducido en la pantalla de inicio del móvil;
  \item Capturas: al menos 2, mejor entre 4 y 6, mostrando las funciones más importantes; acompaña cada una con una frase en lenguaje humano ("registra un gasto en un segundo" es diez veces mejor que "una solución eficiente de contabilidad");
  \item Descripción: la primera frase debe dejar claro "qué es + para quién es"; no acumules calificativos como "el más profesional, el más potente": el posicionamiento en la búsqueda de la tienda depende sobre todo de las palabras clave y de la conversión de descargas; la sinceridad funciona mejor que la exageración.
\end{itemize}

\begin{tcbinfo}{Consejo: compara con los de tu mismo gremio}
Antes de publicar, date una vuelta por las diez primeras apps parecidas a la tuya en la tienda y guarda capturas para comparar: ¿qué estilo tienen sus iconos?, ¿qué puntos fuertes muestran sus capturas?, ¿cómo redactan el título? Imitar no es plagiar: aprende las fórmulas y pon tu propio contenido; es la lección que más rápido arranca cualquier producto nuevo.

\end{tcbinfo}

\booksection{Lista breve de optimización de rendimiento y reporte de errores}

Antes de publicar tu App, estas dos tareas merecen media hora cada una:

\booksubsection{Cuatro puntos de autochequeo de rendimiento}

\begin{enumerate}
  \item Usa LazyColumn para las listas: jamás montes una lista larga con Column; la carga diferida es la primera lección de rendimiento en Android;
  \item No metas las imágenes en resolución original: usa la librería Coil para reescalar y cachear automáticamente; una imagen original de 5 MB basta para atascar un móvil de gama baja;
  \item Evita el trabajo pesado en el hilo principal: base de datos, red y archivos grandes van dentro de corrutinas;
  \item Mide con el Profiler: el analizador de rendimiento incluido en Android Studio muestra memoria, CPU y consumo de batería. Primero mide, después optimiza; no vayas por intuición.
\end{enumerate}

\booksubsection{Reporte de errores}

Que no se caiga en tus pruebas no significa que no se caiga en el móvil del usuario. Antes de publicar, integra Firebase Crashlytics (oficial de Google, gratis): con unas líneas de configuración, cuando la app se caiga se envía automáticamente el stacktrace y en el panel verás “qué modelo, qué versión y en qué línea se cayó”. Publicar sin monitoreo de errores es conducir con los ojos cerrados.

\begin{tcbpractice}{Manos a la obra: "entrega" tu app de gastos una vez}
(1) Añade 3 pruebas unitarias a la lógica, como parseAmount; (2) provoca un fallo tú mismo, atrápalo con Logcat y arréglalo; (3) genera un APK firmado en release y envíalo a un amigo para que lo instale (ese es el verdadero ritual de “obra terminada”); (4) si te sobran fuerzas, crea una cuenta de desarrollador y envía de verdad la app a la tienda. Los tres primeros puntos puedes hacerlos esta misma noche.

\end{tcbpractice}

\begin{tcbsummary}{Resumen del capítulo}
\begin{itemize}
  \item Las tres herramientas de depuración: Logcat para los registros, puntos de interrupción para las variables y Layout Inspector para la interfaz
  \item Leer un fallo: en el stacktrace, de abajo hacia arriba, busca la primera línea con el nombre de tu paquete
  \item Tres niveles de pruebas: unitarias (segundos, aprende estas primero) / instrumentadas / manuales
  \item Para publicar hay que firmar: haz copia del keystore y de la contraseña; si los pierdes, la App "cambia de dueño"
  \item Seis pasos para publicar: registro material subir el AAB formulario de seguridad de datos revisión publicación
\end{itemize}

\end{tcbsummary}
